\chapter{KI-Verwendung}
Im Projekt wurden mehrere KI-Werkzeuge praxisnah eingesetzt. Besonders haeufig kam GitHub Copilot zum Einsatz, sowohl direkt in GitHub als auch in Visual Studio Code. Das hat vor allem bei wiederkehrenden Programmier- und Dokumentationsaufgaben Zeit gespart.

Zusaetzlich wurde Claude Code in einer angepassten Form ("caveman") ausprobiert. Auch Gemini CLI wurde testweise verwendet, um die Ergebnisse und Arbeitsweisen verschiedener Werkzeuge vergleichen zu koennen.

Die GitHub Agents haben sich fuer einfache Tasks als hilfreich erwiesen, zum Beispiel beim Erstellen der grundlegenden Dokumentationsstruktur in diesem Projekt. Gleichzeitig zeigte sich, dass sich die Kosten bei paralleler Nutzung mehrerer KI-Dienste summieren koennen und deshalb aktiv beobachtet werden sollten.
